\documentclass[12pt]{article}
\usepackage[T1]{fontenc}
\usepackage[utf8]{inputenc}
\usepackage{lmodern}
\usepackage{textcomp}
\usepackage{enumitem}
\usepackage{geometry}
\geometry{margin=1in}
\begin{document}
\begin{center}
Answer Key - History - K-12 Level Assessment 25 \
\small (Answers are ordered exactly as in the question paper.)
\end{center}

\section*{Multiple Choice}
\begin{enumerate}[label=\arabic*.]
\item \textbf{Answer:} C
\\ \textit{Options:} A) Caral; B) Stonehenge; C) Poverty Point; D) Gobekli Tepe
\\ \textit{Note:} Poverty Point is the correct answer because it features significant prehistoric earthworks, including concentric ridges arranged in a semicircular pattern, that were constructed around 3,500 years ago, reflecting advanced engineering and social organization of the indigenous cultures at that time.
\item \textbf{Answer:} C
\\ \textit{Options:} A) the height and size of its monumental architecture; B) the amount and frequency of violence; C) the degree of urban planning; D) the amount of surplus wheat that was produced for storage
\\ \textit{Note:} The Indus Valley civilization is considered unmatched in urban planning due to its sophisticated grid layout, advanced drainage systems, and standardized fired-brick buildings, which reflect a high level of organization and foresight not seen in many other early civilizations.
\item \textbf{Answer:} B
\\ \textit{Options:} A) Clovis Complex; B) Denali Complex; C) Folsom Complex; D) Nenana Complex
\\ \textit{Note:} The Denali Complex is the correct answer because it specifically refers to a distinctive lithic technology developed by prehistoric peoples in the Arctic, characterized by the production of wedge-shaped cores, micro-blades, bifacial knives, and burins.
\item \textbf{Answer:} C
\\ \textit{Options:} A) Erectus fossils are found only in Africa.; B) Erectus possessed a smaller brain.; C) Erectus possessed a larger brain.; D) Erectus was a relatively short-lived species.
\\ \textit{Note:} Homo erectus had a larger brain size compared to Homo habilis, indicating a higher level of cognitive ability and more advanced tool-making skills, which marks a significant evolutionary development in early human species.
\item \textbf{Answer:} C
\\ \textit{Options:} A) a willingness to travel great distances; B) a hobby done around the campfire at night; C) increasing awareness and importance of individual identity; D) greater intelligence and willingness to cooperate
\\ \textit{Note:} Randall White argues that the emergence of personal adornment in the Upper Paleolithic period signifies a cultural shift towards recognizing and expressing individual identity, highlighting the increasing importance of personal distinction in social contexts.
\end{enumerate}

\section*{True / False}
\begin{enumerate}[label=\arabic*.]
\setcounter{enumi}{5}
\item \textbf{Answer:} True
\\ \textit{Options:} A) True; B) False
\\ \textit{Note:} Neandertals demonstrated cultural sophistication by creating decorative items such as pendants and necklaces, practicing burial rituals for their deceased, and crafting advanced stone tools, indicating their capability for symbolic thought and social behaviors.
\item \textbf{Answer:} False
\\ \textit{Options:} A) True; B) False
\\ \textit{Note:} Radiocarbon dating is based on the decay of carbon-14 in organic materials, not on tree-ring growth patterns, which is a separate method known as dendrochronology.
\item \textbf{Answer:} True
\\ \textit{Options:} A) True; B) False
\\ \textit{Note:} Brushes, trowels, and dental picks are all specialized tools used by archaeologists to carefully excavate and remove soil around artifacts and features without damaging them.
\item \textbf{Answer:} False
\\ \textit{Options:} A) True; B) False
\\ \textit{Note:} Artifacts found in a landfill represent a secondary context for archaeological study because they have been disturbed and relocated from their original site of use, making it difficult to ascertain their historical significance and context.
\item \textbf{Answer:} True
\\ \textit{Options:} A) True; B) False
\\ \textit{Note:} The statement is true because post-Pleistocene Europe saw a variety of ecosystems that supported a wide range of flora and fauna, leading to diverse food sources such as wild game, fish, nuts, fruits, and edible plants available for human consumption.
\end{enumerate}

\section*{Long Form Response}
\begin{enumerate}[label=\arabic*.]
\setcounter{enumi}{10}
\item \textbf{Answer:} Following the classical period around A.D. 800, the Mayan civilization experienced significant changes, including a decline in many of its southern cities. Factors contributing to this decline included environmental challenges, such as drought, overpopulation, and warfare between city-states. As a result, many Mayans moved their centers northward to areas like the Yucatán Peninsula. This shift allowed for a resurgence of Mayan culture and power, particularly during the Postclassic period. However, this resurgence was not permanent, and the civilization ultimately fell into decline once again due to ongoing environmental and social pressures.
\\ \textit{Note:} correct because it accurately identifies the decline of southern Mayan cities due to environmental factors and warfare, and it highlights the subsequent migration and cultural resurgence in the northern regions during the Postclassic period.
\item \textbf{Answer:} Cranial capacity in Homo erectus, which averaged just under 1000 cc, is significant as it suggests an increase in brain size compared to earlier hominins. This larger brain size likely contributed to enhanced cognitive abilities, allowing Homo erectus to develop more complex tools, engage in problem-solving, and adapt to diverse environments. Additionally, these cognitive advancements may have influenced their social behaviors, enabling them to communicate more effectively and form larger social groups. As a result, Homo erectus could better cooperate in hunting and gathering, demonstrating an evolution in their social structure and interactions during prehistoric times.
\\ \textit{Note:} correct because it connects the increased cranial capacity of Homo erectus to enhanced cognitive abilities and social behaviors, demonstrating how this physical trait likely facilitated tool development, problem-solving, and effective communication, which were crucial for their survival and adaptation in prehistoric environments.
\end{enumerate}

\vfill
\noindent\textit{End of Answer Key}
\end{document}